\section*{الفصل \ref{ch:b1:populations} --- الجمهرات والديموغرافيا}

\begin{solution}{exo:b1:populations:1}
\omterm{def:b1:populations:population}{الجمهرة}: أفراد نوع واحد في مكان وزمن واحدين، يستطيعون التزاوج فيما
بينهم. والكثافة: الأفراد في وحدة المساحة أو الحجم. والتوزع: نمطهم
المكاني --- متكتل (قطيع، أو بقعة قُرّاص)، ومنتظم (أطيش تعشش على بعد
نقرة، وجنبات الكريوزوت)، وعشوائي (هندباء برية على مرج).
\end{solution}

\begin{solution}{exo:b1:populations:2}
$l_x$: كسر الفوج الحي عند العمر $x$. و$m_x$: البنات لكل
أنثى عمرها $x$. و$R_0 = \sum l_x m_x$: البنات لكل أنثى مولودة
على مدى حياتها. و$T = \sum x l_x m_x / R_0$: متوسط عمر الأمهات.
و$R_0 = 1$: كل أنثى تحلّ محل نفسها؛ \omterm{def:b1:populations:population}{فالجمهرة} مستقرة.
\end{solution}

\begin{solution}{exo:b1:populations:3}
النمط III: نحو 15 من 1000، أي \qty{1.5}{\%}. والنمط I: نحو 700 من
1000، أي \qty{70}{\%}.
\end{solution}

\begin{solution}{exo:b1:populations:4}
$\dd N/\dd t = rN$ و$\dd N/\dd t = rN(1 - N/K)$. والمقدار $r$ هو
معدل النمو للفرد حين لا تكون الموارد محدودة (الولادات ناقص
الوفيات لكل فرد في وحدة الزمن)؛ و$K$ هي \omterm{thm:b1:populations:logistic}{السعة الحمولية}، أي
الحجم الذي يتوقف عنده النمو.
\end{solution}

\begin{solution}{exo:b1:populations:5}
$r = \ln(1000)/5 = \qty{1.38}{h^{-1}}$؛ والتضاعف $\ln 2/1.38 =
\qty{0.5}{h}$.
\end{solution}

\begin{solution}{exo:b1:populations:6}
$l_x m_x = 0, 0.5, 0.6, 0.2, 0$: أي $R_0 = 1.3$. و$T = (1\times 0.5 +
2\times 0.6 + 3\times 0.2)/1.3 = 2.3/1.3 = 1.77$. و$r = \ln 1.3/1.77 =
\qty{0.15}{}$ في وحدة الزمن: فهي تنمو.
\end{solution}

\begin{solution}{exo:b1:populations:7}
$N = 50\times 70/7 = 500$. ولم يبق موسومًا إلا 40 سمكة: فيكون $N = 40\times
70/7 = 400$؛ والتقدير الأول كان أعلى من الحقيقة بالربع (إذ كان كسر
الموسوم من \omterm{def:b1:populations:population}{الجمهرة} أصغر مما فُرض).
\end{solution}

\begin{solution}{exo:b1:populations:8}
$0.2\times 100\times 0.9 = 18$؛ و$0.2\times 500\times 0.5 = 50$؛
و$0.2\times 900\times 0.1 = 18$ في السنة. والأعظم $rK/4 = 50$ في السنة، عند
$N = 500$.
\end{solution}

\begin{solution}{exo:b1:populations:9}
الهندباء البرية: $r$ (آلاف البذور، حولية). والفيل: $K$ (عجل واحد
كل أربع سنوات، عمر طويل، رعاية). وسمكة القدّ: $r$ (ملايين البيوض، بلا رعاية).
والغوريلا: $K$ (رضيع واحد كل أربع سنوات، عقود من العمر). وذبابة المنزل: $r$
(مئات البيوض، تنضج في عشرة أيام).
\end{solution}

\begin{solution}{exo:b1:populations:10}
$r = \ln(29/10)/2 = \qty{0.53}{h^{-1}}$؛ و$K \approx 660$. وعند \qty{8}{h}:
$660/(1 + 65e^{-4.24}) = 660/(1 + 0.94) = 340$، في مقابل 351 معدودة.
والأسرع عند $K/2 = 330$، وتُبلَغ عند $t = \ln 65/0.53 = \qty{7.9}{h}$،
حيث $\dd N/\dd t = rK/4 = 87\times 10^5$ \omterm{def:b1:cell-unit-of-life:cell}{خلية} في الميليلتر في
الساعة.
\end{solution}

\begin{solution}{exo:b1:populations:11}
يقتضي التنظيم أن يهبط المعدل للفرد مع ارتفاع $N$، حتى يعود
$N$ نحو توازن بعد اضطراب؛ أما عامل يقتل كسرًا ثابتًا فيغيّر
$N$ ولا يغيّر تعلق المعدل بالحجم $N$، فتستأنف \omterm{def:b1:populations:population}{الجمهرة} بعده النمو نفسه.
والطقس يحدّ الحشرات بأن يضربها مرارًا، وعشوائيًا، فتُرد \omterm{def:b1:populations:population}{الجمهرة} إلى
الوراء قبل أن تقارب أي سقف: أي \omterm{def:b1:populations:population}{جمهرة} غير منظَّمة يبقيها الاضطراب
المتكرر منخفضة، ويتوقف حجمها الوسطي على تواتر المواسم السيئة لا
على \omterm{thm:b1:populations:logistic}{سعة حمولية}.
\end{solution}

\begin{solution}{exo:b1:populations:12}
جمهرتان بالحجم نفسه، إحداهما \omterm{def:b1:nucleic-acids:nucleotide}{بقاعدة} فتية عريضة والأخرى بفئات
عمرية متساوية، لهما مستقبلان مختلفان: فالأولى ستنمو عقودًا حتى لو
هبط $m_x$ فيها إلى الإحلال الآن، لأن صغارها الكثيرين لم يتكاثروا
بعد (وهو الزخم الديموغرافي)؛ والثانية لن تنمو. والعدد $N$ يسجل الماضي؛
أما التوزيع العمري وجدول $m_x$ عليه فهما ما تعمل عليه المعادلات،
والهرم يتنبأ بحجم الجيل التالي بينما لا يقول $N$ عنه شيئًا.
\end{solution}

\begin{solution}{pb:b1:populations:1}
\textbf{1.} $r = \ln(29/10)/2 = \qty{0.53}{h^{-1}}$.
\textbf{2.} $\ln 2/0.53 = \qty{1.3}{h}$.
\textbf{3.} $K \approx 660\times 10^5 = \qty{6.6e7}{}$ \omterm{def:b1:cell-unit-of-life:cell}{خلية} في
الميليلتر.
\textbf{4.} $N(t) = 660/(1 + 65e^{-0.53t})$: فعند 4 ساعات، $660/(1 + 7.8) =
75$ (والمعدود 71)؛ وعند 8 ساعات، 340 (والمعدود 351)؛ وعند 12 ساعة، $660/(1 + 0.112) = 594$
(والمعدود 595).
\textbf{5.} عند $K/2 = 330$، وتُبلَغ عند $t = \ln 65/0.53 = \qty{7.9}{h}$؛
والمعدل الأعظمي $rK/4 = 0.53\times 660/4 = 87\times 10^5$ \omterm{def:b1:cell-unit-of-life:cell}{خلية} في
الميليلتر في الساعة.
\textbf{6.} $10e^{0.53\times 12} = 10\times 578 = 5780$: أي عشرة أمثال
595 الملاحظة.
\textbf{7.} نفاد السكر (أو الأكسجين)؛ وتراكم الإيثانول والأحماض
التي تسمم \omterm{def:b1:cell-unit-of-life:cell}{الخلايا}.
\textbf{8.} $q_0 = 0.40$؛ و$q_4 = 1 - 0.34/0.40 = 0.15$: أي وفيات
مبكرة ثقيلة ثم معدل أثبت --- بين النمط II والنمط III، كأكثر
الثدييات الكبيرة في البرية.
\textbf{9.} $l_x m_x = 0, 0, 0.312, 0.368, 0.320, 0.272, 0.182, 0.080,
0.018, 0$: أي $R_0 = 1.55$.
\textbf{10.} $\sum x l_x m_x = 0.624 + 1.104 + 1.280 + 1.360 + 1.092 +
0.560 + 0.144 = 6.16$؛ و$T = 6.16/1.55 = \qty{4.0}{yr}$.
\textbf{11.} $r = \ln 1.55/4.0 = \qty{0.11}{yr^{-1}}$؛ والتضاعف في
\qty{6.3}{yr}.
\textbf{12.} $200e^{1.1} = 600$ أيّل.
\textbf{13.} $N(10) = 600/(1 + 2e^{-1.1}) = 600/1.67 = 360$؛ ويتجاوز 500
حين $(K - N_0)/N_0\,e^{-rt} = 0.2$، أي $e^{-rt} = 0.1$، فيكون $t = \ln
10/0.11 = 21$ سنة.
\textbf{14.} $120\times 150/18 = 1000$ سمكة فرخ.
\textbf{15.} 100 موسومة: $100\times 150/18 = 833$؛ فالتقدير الأول
كان أعلى من الحقيقة (إذ كان السمك الموسوم الطليق أقل مما فُرض).
\textbf{16.} $rK/4 = 100$ في السنة، عند $N = 500$.
\textbf{17.} $T_2 = \ln 2/0.4 = \qty{1.7}{yr}$.
\textbf{18.} الصيد $hN$ عند $h = 0.2$: فالمعادلة $rN(1 - N/K) = hN$ تعطي $N^* =
K(1 - h/r) = 1000\times 0.5 = 500$ وصيدًا قدره \num{100} في السنة ---
أي دون أعظم غلة مستدامة، لكن الكسر يضبط نفسه على الرصيد: فإذا
انتصف الرصيد انتصف الصيد، وعادت \omterm{def:b1:populations:population}{الجمهرة}.
\textbf{19.} عند 500: $0.4\times 500\times 0.5 = 100 < 120$: فتتراجع
\omterm{def:b1:populations:population}{الجمهرة}. وعند 300: $0.4\times 300\times 0.7 = 84 < 120$:
فتتراجع أسرع --- فمتى نزلت دون $K/2$ ساق حصادٌ ثابت فوق الغلة
العظمى الجمهرةَ إلى الانقراض.
\textbf{20.} عند $K/2$ يكون الفائض في ذروته لكن أي خطأ ---
\omterm{def:b1:populations:population}{جمهرة} قُدّرت أعلى بالخمس، أو سنة سيئة تقصّ $r$ ---
يحوّل أخذًا مستدامًا إلى تراجع، ودون $K/2$ يتقلص الفائض مع تقلص
\omterm{def:b1:populations:population}{الجمهرة}: فالتوازن غير مستقر أمام الإفراط في الحصاد. والحصاد فوق
$K/2$ قليلًا يترك هامشًا.
\textbf{21.} ينتصف $R_0$ إلى $0.78 < 1$: فيتقلص القطيع حتى يهبط
دون 400 ويتعافى $m_x$ --- وهو عامل متعلق بالكثافة، ينظّم القطيع
قرب 400.
\textbf{22.} لا: فهو يزيل \qty{40}{\%} مهما يكن الحجم ويترك
المعدل للفرد كما هو؛ وفي السنوات التالية ينمو القطيع ثانية بالمعدل
$r$ نفسه من بداية أدنى --- أي اضطراب لا تنظيم.
\textbf{23.} $0.11N(1 - N/600) = 30$: أي $N^2 - 600N + 163\,600 = 0$،
ومنه $N = 300\pm\sqrt{90\,000 - 163\,600}$ --- ولا حل حقيقي: فالعدد 30 في السنة
تتجاوز الفائض الأعظم $rK/4 = 16.5$، وتسوق الذئاب
القطيع إلى الانقراض. ومع 15 في السنة: $N = 300\pm\sqrt{90\,000 -
81\,800} = 300\pm 91$، أي 209 و391؛ والأعلى مستقر
(ففوقه يكون الفائض أقل من 15 فيعود القطيع نازلًا، ودونه أكثر
فيرتفع)، والأدنى غير مستقر.
\textbf{24.} تستجيب أعداد الوشق للغذاء بتأخر: فارتفاع الأرانب يرفع
ولادات الوشق وبقياه على مدى السنة أو السنتين التاليتين، وانهيار
الأرانب لا يجيع الوشق إلا بعد نفاد شحمه وأشبال الموسم.
\textbf{25.} الخميرة: $K = \qty{6.6e7}{}$ \omterm{def:b1:cell-unit-of-life:cell}{خلية} في الميليلتر، و$r =
\qty{0.53}{h^{-1}}$. والأيائل: $R_0 = 1.55$، و$T = \qty{4.0}{yr}$، و$r =
\qty{0.11}{yr^{-1}}$.
\end{solution}
